% GENERATED FILE. Edit canonical lessons, then run npm run generate:books.
% canonical-source-hash: fnv1a64:9fa14e62e4eacbde
% canonical-lessons: RU-C144-boyatsya, RU-C144-nravitsya, RU-C144-nenavidet, RU-C144-khvalit, RU-C144-rugat

\chapter{To be afraid, To please, To hate, To praise, To scold}
\label{ch:ru-to-be-afraid-to-please-to-hate-to-praise-to-scold}
\hlchaptermodality{144}

\begin{chapteropening}
\textbf{By the end of this chapter:} \emph{I can say to be afraid, to please, to hate, to praise and to scold.}

Complete the last lesson of chapter 144: I can say to be afraid, to please, to hate, to praise and to scold.
\end{chapteropening}

\section[\texorpdfstring{\ru{бояться} (boyátsya)}{бояться (boyátsya)}]{\ru{бояться} (boyátsya) — to be afraid}

\label{lesson:RU-C144-boyatsya}



\begin{quote}
Before the new one: say the Russian for to come, then the Russian for to leave.
\end{quote}

\subsection*{You'll want to know: \ru{бояться}}
\textbf{\ru{бояться}} (\emph{boyátsya}) — \textquotedblleft{}to be afraid\textquotedblright{}.

\begin{grammarlens}[title={What you've built}]
One more word for this chapter.
\end{grammarlens}

\subsection*{Your turn}
\begin{itemize}
\raggedright
  \item \emph{Say it:} \emph{boyátsya}
  \item \emph{Say it:} \emph{boyátsya}, once more
  \item \emph{Say it:} say \emph{boyátsya}
  \item \emph{Recall:} say the Russian for to come, then the Russian for to leave, then say \emph{boyátsya} again
\end{itemize}

\begin{tcolorbox}[breakable,colback=teal!4,colframe=teal!35!black,title={Before you move on}]
What does \ru{бояться} mean? (\textquotedblleft{}To be afraid\textquotedblright{}.) Say it once more.
\end{tcolorbox}
\section[\texorpdfstring{\ru{нравиться} (nrávitsya)}{нравиться (nrávitsya)}]{\ru{нравиться} (nrávitsya) — to please, to be liked}

\label{lesson:RU-C144-nravitsya}



\begin{quote}
Before the new one: say the Russian for to leave, then the Russian for to be afraid.
\end{quote}

\subsection*{You'll want to know: \ru{нравиться}}
\textbf{\ru{нравиться}} (\emph{nrávitsya}) — \textquotedblleft{}to please, to be liked\textquotedblright{}.

\begin{grammarlens}[title={What you've built}]
One more word for this chapter.
\end{grammarlens}

\subsection*{Your turn}
\begin{itemize}
\raggedright
  \item \emph{Say it:} \emph{nrávitsya}
  \item \emph{Say it:} \emph{nrávitsya}, once more
  \item \emph{Say it:} say \emph{nrávitsya}
  \item \emph{Recall:} say the Russian for to leave, then the Russian for to be afraid, then say \emph{nrávitsya} again
\end{itemize}

\begin{tcolorbox}[breakable,colback=teal!4,colframe=teal!35!black,title={Before you move on}]
What does \ru{нравиться} mean? (\textquotedblleft{}To please, to be liked\textquotedblright{}.) Say it once more.
\end{tcolorbox}
\section[\texorpdfstring{\ru{ненавидеть} (nenavídet)}{ненавидеть (nenavídet)}]{\ru{ненавидеть} (nenavídet) — to hate}

\label{lesson:RU-C144-nenavidet}



\begin{quote}
Before the new one: say the Russian for to be afraid, then the Russian for to please.
\end{quote}

\subsection*{You'll want to know: \ru{ненавидеть}}
\textbf{\ru{ненавидеть}} (\emph{nenavídet}) — \textquotedblleft{}to hate\textquotedblright{}.

\begin{grammarlens}[title={What you've built}]
One more word for this chapter.
\end{grammarlens}

\subsection*{Your turn}
\begin{itemize}
\raggedright
  \item \emph{Say it:} \emph{nenavídet}
  \item \emph{Say it:} \emph{nenavídet}, once more
  \item \emph{Say it:} say \emph{nenavídet}
  \item \emph{Recall:} say the Russian for to be afraid, then the Russian for to please, then say \emph{nenavídet} again
\end{itemize}

\begin{tcolorbox}[breakable,colback=teal!4,colframe=teal!35!black,title={Before you move on}]
What does \ru{ненавидеть} mean? (\textquotedblleft{}To hate\textquotedblright{}.) Say it once more.
\end{tcolorbox}
\section[\texorpdfstring{\ru{хвалить} (khvalít)}{хвалить (khvalít)}]{\ru{хвалить} (khvalít) — to praise}

\label{lesson:RU-C144-khvalit}



\begin{quote}
Before the new one: say the Russian for to please, then the Russian for to hate.
\end{quote}

\subsection*{You'll want to know: \ru{хвалить}}
\textbf{\ru{хвалить}} (\emph{khvalít}) — \textquotedblleft{}to praise\textquotedblright{}.

\begin{grammarlens}[title={What you've built}]
One more word for this chapter.
\end{grammarlens}

\subsection*{Your turn}
\begin{itemize}
\raggedright
  \item \emph{Say it:} \emph{khvalít}
  \item \emph{Say it:} \emph{khvalít}, once more
  \item \emph{Say it:} say \emph{khvalít}
  \item \emph{Recall:} say the Russian for to please, then the Russian for to hate, then say \emph{khvalít} again
\end{itemize}

\begin{tcolorbox}[breakable,colback=teal!4,colframe=teal!35!black,title={Before you move on}]
What does \ru{хвалить} mean? (\textquotedblleft{}To praise\textquotedblright{}.) Say it once more.
\end{tcolorbox}
\section[\texorpdfstring{\ru{ругать} (rugát)}{ругать (rugát)}]{\ru{ругать} (rugát) — to scold}

\label{lesson:RU-C144-rugat}



\begin{quote}
Before the new one: say the Russian for to hate, then the Russian for to praise.
\end{quote}

\subsection*{You'll want to know: \ru{ругать}}
\textbf{\ru{ругать}} (\emph{rugát}) — \textquotedblleft{}to scold\textquotedblright{}.

\begin{grammarlens}[title={What you've built}]
One more word for this chapter.
\end{grammarlens}

\subsection*{Your turn}
\begin{itemize}
\raggedright
  \item \emph{Say it:} \emph{rugát}
  \item \emph{Say it:} \emph{rugát}, once more
  \item \emph{Say it:} say \emph{rugát}
  \item \emph{Recall:} say the Russian for to hate, then the Russian for to praise, then say \emph{rugát} again
\end{itemize}

\begin{tcolorbox}[breakable,colback=teal!4,colframe=teal!35!black,title={Before you move on}]
What does \ru{ругать} mean? (\textquotedblleft{}To scold\textquotedblright{}.) Say it once more.
\end{tcolorbox}
